\documentclass[11pt, a4paper]{article}
\usepackage[inner=1.5cm,outer=1.5cm,top=2.5cm,bottom=2.5cm]{geometry}
\usepackage{graphicx}
\usepackage{fancyhdr, lastpage, pmboxdraw}
\usepackage[usenames,dvipsnames]{color}
\definecolor{darkblue}{rgb}{0,0,.6}
\definecolor{darkred}{rgb}{.7,0,0}
\definecolor{darkgreen}{rgb}{0,.6,0}
\usepackage[colorlinks,pdfusetitle,urlcolor=darkblue,citecolor=darkblue,linkcolor=darkred,bookmarksnumbered,plainpages=false]{hyperref}
\usepackage{listings}
\usepackage{caption}
\usepackage{tabularx}
\usepackage{verbatim}
\usepackage{fancyvrb}

\newcommand{\custom}[1]{{\color{blue} #1}}
\definecolor{OliveGreen}{cmyk}{0.64,0,0.95,0.40}
\definecolor{lightlightgray}{gray}{0.93}

\pagestyle{fancyplain}
\fancyhf{}
\lhead{ \fancyplain{}{Statistical Programming in R} }
\rhead{ \fancyplain{}{DATA 412/612} }
\fancyfoot[RO, LE] {page \thepage\ of \pageref{LastPage} }
\thispagestyle{plain}

\lstset{
language=R,
basicstyle=\ttfamily\footnotesize,
keywordstyle=\color{OliveGreen},
commentstyle=\color{gray},
numbers=left,
numberstyle=\tiny,
stepnumber=1,
numbersep=5pt,
backgroundcolor=\color{lightlightgray},
frame=none,
tabsize=2,
captionpos=t,
breaklines=true,
breakatwhitespace=false,
showspaces=false,
showtabs=false,
columns=flexible
}

\begin{document}
\begin{center}
{\Large \textsc{\textbf{A Short Guide to Exploratory Data Analysis}}}
\\ \textbf{Statistical Programming in R}\\ \textbf{DATA 412/612} \\
Department of Mathematics and Statistics
\\
American University
\end{center}
\begin{center}
Fall 2026
\end{center}

\section*{What EDA is}

Exploratory data analysis is how you get to know a table: ask a question, look at the data, then ask a better question. It is not a dump of every \texttt{ggplot} you tried, and it is not a final confirmatory study. You are allowed to be surprised.

\emph{R for Data Science} puts it this way: two questions come up over and over.
\begin{itemize}
    \item What kind of variation is there \emph{within} a variable?
    \item What kind of covariation is there \emph{between} variables?
\end{itemize}
That is why the project asks for real data with at least three numeric columns, plus categories, dates, and some text.

\section*{A reasonable order}

You will jump backward when a plot shows a cleaning problem. That is normal. For the write-up, try this order.

\begin{enumerate}
    \item \textbf{One row.} Read the codebook. What is the observational unit?
    \item \textbf{Import.} \texttt{readr} / \texttt{readxl} / \texttt{haven}, then \texttt{glimpse()} and a few \texttt{count()}s. Fix types (numbers or dates stored as text are common).
    \item \textbf{Tidy and clean.} Rename, parse dates, recode labels, join or pivot, drop impossible values. Write down what you dropped.
    \item \textbf{Missingness.} How much, and does it depend on year or group?
    \item \textbf{One variable at a time.} Histograms or boxplots for numeric variables; bar charts for categories (reorder the levels); counts over time for dates.
    \item \textbf{Two or three variables.} Scatterplots, boxplots by group, line plots over time. Facet or color for a third variable. If two numeric variables do not line up in a straight cloud, see below.
    \item \textbf{Odd points.} Outliers, zeros, and codes like 999 are part of the data. See if your pattern still holds without them.
    \item \textbf{New questions.} The data-driven hypotheses should not be a copy of the initial ones.
    \item \textbf{A story.} Keep the plots that you can explain. The rest can go in an appendix or be deleted. The knitted PDF is at most 20 pages.
\end{enumerate}

\section*{Type $\rightarrow$ plot $\rightarrow$ package}

\begin{center}
\renewcommand{\arraystretch}{1.2}
\begin{tabularx}{\textwidth}{>{\raggedright\arraybackslash}p{1.45in} X >{\raggedright\arraybackslash}p{1.65in}}
\hline
\textbf{You have} & \textbf{Look at} & \textbf{Tools} \\
\hline
one numeric variable & histogram, boxplot, quantiles & \texttt{ggplot2}, \texttt{summarise} \\
numeric vs.\ numeric & scatterplot first; smooth / log / bins if it is not a line & \texttt{geom\_point}, \texttt{geom\_smooth} \\
numeric vs.\ category & boxplots, grouped summaries & \texttt{group\_by}, \texttt{forcats} \\
category vs.\ category & counts, bar charts of proportions & \texttt{count}, \texttt{geom\_bar} \\
dates & parse, then counts or means over time & \texttt{lubridate} \\
text & flags, splits, cleaned labels, then plot & \texttt{stringr} \\
many files / groups & the same pipeline repeatedly & \texttt{purrr} \\
related tables or wide data & joins or \texttt{pivot\_longer} & \texttt{dplyr}, \texttt{tidyr} \\
\hline
\end{tabularx}
\end{center}

Reorder factors with \texttt{fct\_reorder()}, lump rare levels with \texttt{fct\_lump()}, and put units on the axes. Captions should say something, not ``Plot 3.''

\section*{When two numeric variables are not a straight line}

Pearson correlation and a fitted straight line only describe a \emph{linear} association. They can be near zero when there is a clear curve, and they can look large when the cloud is just two groups sitting apart. Always start with a scatterplot.

\begin{itemize}
    \item Add a smooth: \texttt{geom\_smooth()} (loess for smaller $n$; for large $n$ use \texttt{method = "gam"} or bin first). Read the curve as a description of the cloud, not as a causal model.
    \item Try a log scale on $x$ or $y$ when the pattern bends or the spread grows with the mean. If the plot is closer to a line on that scale, summarize there and say so.
    \item Cut one variable into bins (\texttt{cut()} or \texttt{ntile()}) and plot the mean or median of the other against the bin. That is a simple check that does not assume a line.
    \item Spearman correlation is for a monotone pattern (always up, or always down) that is not linear. Report it \emph{with} the plot, not instead of the plot.
    \item If you fit \texttt{lm()}, look at residuals vs.\ fitted. A curve there means the line missed the relationship.
    \item Facet or color by a categorical variable. A bend in the pooled scatter can be two groups with different intercepts; a straight pooled cloud can hide opposite slopes.
\end{itemize}

Do this before you write that two variables are related or that they are not. Overplotting can hide a curve; use \texttt{alpha} or \texttt{geom\_hex()} when $n$ is large. The smooth is still exploratory: say what you see, how large it is, and where it might be an artifact of missingness, outliers, or a third variable.

\section*{A small pipeline}

\begin{lstlisting}
library(tidyverse)
library(lubridate)

raw <- read_csv("../data/listings.csv")

clean <- raw %>%
  mutate(
    price = parse_number(price),
    last_review = ymd(last_review),
    neighborhood = fct_lump(neighborhood, n = 10),
    has_parking = str_detect(amenities, "parking")
  ) %>%
  filter(price > 0, price < 2000)

clean %>%
  group_by(neighborhood) %>%
  summarise(n = n(), median_price = median(price, na.rm = TRUE),
            .groups = "drop") %>%
  mutate(neighborhood = fct_reorder(neighborhood, median_price)) %>%
  ggplot(aes(median_price, neighborhood)) +
  geom_col() +
  labs(x = "Median nightly price (USD)", y = NULL)
\end{lstlisting}

Several files of the same shape:

\begin{lstlisting}
files <- list.files("../data", pattern = "\\.csv$", full.names = TRUE)
trips <- files %>% map_dfr(read_csv, .id = "source")
\end{lstlisting}

A nonlinear numeric--numeric plot:

\begin{lstlisting}
clean %>%
  ggplot(aes(reviews, price)) +
  geom_point(alpha = 0.3) +
  geom_smooth(se = FALSE) +
  scale_x_log10() +
  facet_wrap(~ room_type) +
  labs(x = "Number of reviews", y = "Nightly price (USD)")
\end{lstlisting}

\section*{Writing it up}

For each important figure: one sentence of setup, the figure, then two or three sentences that state the pattern, its size, and a caveat. Connect figures instead of listing them. Do not claim a relationship that the plot does not show, and do not treat a smooth or a correlation as the end of the argument.

A sequence that usually works: structure and missingness; the main numeric variables; the main groups; the relationship you care about (stratified); the same thing over time if you have dates; one surprise.

The EDA chapter in \emph{R for Data Science} is the main reading: \url{https://r4ds.hadley.nz/eda}.

\end{document}
